% Original PDF pp.6--10: dedication and complete preface.
\clearpage
\thispagestyle{empty}
\vspace*{.25\textheight}
\begin{center}\large Kariへ\end{center}
\clearpage
\sectionlink*{序文}

過去10年のあいだに、デコヒーレンスは、量子物理学の基礎理論からナノ工学への応用まで、あらゆる種類の研究で使われる、至るところに見られる科学用語になった。デコヒーレンスは、量子力学の誕生にまでさかのぼる、長年の基礎的問題の解決として称賛されてきた。また、量子コンピューター（これもまた、どこにでもあるように見える研究分野である）のような次世代技術に対する主要な障害として、忌み嫌われてもきた。そして、デコヒーレンスはさまざまな実験で直接観測されているにもかかわらず、その適用範囲と意味は、しばしば誤解され、誤って伝えられてきた。デコヒーレンスは、解釈に関する哲学的な問いから、実験や工学的応用における実際的な問題に至るまで、物理学のさまざまな側面に関わるため、実に魅力的な研究主題である。

本書は、デコヒーレンスのこうしたさまざまな魅力的側面を、理解しやすく自己完結した形で読者に紹介する。とくに焦点を当てるのは、デコヒーレンスと、いわゆる量子から古典への移行との関係である。すなわち、私たちの周囲の巨視的世界が古典的に見えることが、基礎にある量子的基盤から現れるのを、デコヒーレンスはどのように説明できるのか、という問いである。

本書は、理論と実験の両領域におけるデコヒーレンスの多くの側面に読者が親しめるよう、比較的広い範囲を扱う。本書全体を通じて、一方にあるデコヒーレンス研究の概念的な考え方と、他方にある形式的・数学的な詳細とのあいだに、適切な均衡を保つよう努めた。本書は、デコヒーレンスを純粋な量子現象として正しく理解するための基盤を築くとともに、その帰結と成果を正しく解釈することの重要性を強調する。

デコヒーレンスを学ぶことのすばらしい点の一つは、その含意と応用がいかに広大であっても、基本的な考え方と形式構造は、実は非常に明快で単純だということである。原則として、可能なかぎり、重要な一般的洞察を複雑な数学的計算でわかりにくくすることを避ける。本書の説明と導出は、すべてとはいわなくても大部分について、量子力学の形式体系の基礎知識があれば理解できるはずである。いくつかの節には、やむをえず、かなり長い数学的考察が含まれる（第4章と第5章のマスター方程式の導出が、おそらく最も顕著な例である）。しかし、デコヒーレンス研究の基礎にあるこうした形式構造にあまり関心のない読者は、それらの節をざっと見るだけでも、あるいはまったく飛ばしても、他の部分の理解を大きく損なわずに済むはずである。同時に、本書に含まれるより進んだ内容は、デコヒーレンスについてすでに多少の知識があり、自己完結した詳細な参考書を求めている、研究に携わる物理学者にも役立つだろう。量子力学の基礎に関心をもつ物理学の哲学者も、本書の至るところ、とくに第1、2、8、9章に、多くの興味深い内容を見いだすはずである。

本書の構成は以下のとおりである。第1章では、基本的な考え方と概念のいくつかを紹介し、デコヒーレンスを初めて「鳥瞰」する。量子論の長年の問題のいくつかを扱うには、「開いた」量子系を考えることが重要であると強調し、この認識とデコヒーレンス研究の考え方が根付くまでに、なぜ半世紀以上も要したのかを考える。

本書の中心は第2章であり、そこではデコヒーレンスの主要な概念と形式的記述を導入し、詳しく論じる。まず、量子状態（および古典状態との違い）、重ね合わせの原理、量子もつれ、密度行列といった、量子力学の基本概念を分析する。これらを正しく把握することは、デコヒーレンスを確実に理解するうえで、非常に重要であることが明らかになる。次に、「量子測定問題」として知られるようになった問題の、さまざまな構成要素を例示して論じる。この問題には、量子力学の形式体系の物理的意味と、この形式体系が知覚される周囲の世界とどう関係するかについて、今日まで共通の理解に達することを妨げてきた、多くの基本的な概念上の難点が凝縮されている。デコヒーレンスはこの問題の各構成要素に直接の含意をもつため、測定問題はデコヒーレンスとも密接に関係している。

続いて、よく知られた二重スリット実験の文脈で、デコヒーレンスの基本概念を説明する。この方法により、読者は、環境による「監視」と量子もつれの帰結としてのデコヒーレンスを、かなり自然に理解できるだろう。また、Bohrの有名な「相補性原理」についての現代的な見方も得られる。系と環境のもつれ、および縮約密度行列によってデコヒーレンスを定式化し、その二つの主要な帰結、すなわち環境による量子干渉の抑制と、環境との相互作用による優先「ポインター」状態の選択を論じる。

このようにしてデコヒーレンスの考え方と形式体系に慣れた後は、後続の章を順番に読んでもよいし、とくに関心のある章に集中してもよい。各章は、デコヒーレンスの特定の側面について、かなり自己完結した議論を提供するように構成されている。

第3章では、光子や空気分子のような環境粒子との衝突によって量子的対象に生じるデコヒーレンスを記述する、非常に重要な模型を考える。この散乱誘起デコヒーレンスは自然界に遍在し、日常世界の巨視的スケールにおける量子から古典への移行を記述するうえで、きわめて重要である。

次の第4章では、物理的に関心のある多くの場合について、デコヒーレンス模型の力学を決定する一般的方法を与える、マスター方程式の形式体系を導入する。かなり直接的で直観的な仕方で多くのデコヒーレンスの問題を扱えるようにする、重要なBorn--Markovマスター方程式の導出に、ある程度の時間をかける。

第5章では、系と環境の模型の広い一群を、少数の「標準的な」デコヒーレンス模型へと帰着できることを示す。続いて、これらの模型を詳しく分析する。とくに、位相空間における、よく知られた古典的なNewton力学的軌道への量子的近似とみなせる、いわゆる量子Brown運動を論じる。また、量子計算の文脈で近年さらに注目されている、有名なスピン・ボソン模型も導入する。

これだけ多くの理論的内容の後では、読者もきっと一休みしたくなるだろう。そこで第6章では、実験室でデコヒーレンスの漸進的な作用、したがって量子世界から古典的領域への移行を直接観測することを可能にした、魅力的な実験のいくつかを述べる。

第7章では、いくらか話題を変え、過去10年間に大きな関心を集めてきた量子計算の分野に入る。この分野でデコヒーレンスが果たす決定的な役割を説明する。次に、量子誤り訂正、デコヒーレンスのない部分空間、環境工学といった、洗練されてはいるが最終的には理解しやすい方法によって、デコヒーレンスの影響をどう軽減できるかを述べる。

第8章では、量子力学の主要な解釈のいくつかに対する、デコヒーレンスの含意を論じる。正統的解釈とCopenhagen解釈、相対状態解釈、物理的収縮模型、様相解釈、Bohm力学など、最も一般的な解釈を、デコヒーレンスがどのように強化し、再定義し、あるいはそれらに疑問を投げかけるのかを述べる。

最後の第9章では、量子論における観測者の役割と、脳におけるデコヒーレンス過程の問題を論じる。まず、なぜこの問いに関心が向けられるのかを説明し、次に、脳内でのデコヒーレンスの効率のよさを示す、具体的な模型計算のいくつかを概観する。これらの結果の含意を、とくに観測者に基づく測定問題の「主観的」な解決との関係で論じる。

記法について、ひと言述べておく。明示的な数値評価が重要な状況を除き、本書の大部分では$\hbar\equiv1$とした。この方法によって、必要なときに読者が数値を導き、再現する能力を損なわずに、記法をできるだけ明快に保てたことを願っている。

本書を可能にするために貢献してくださった方々は数多い。まず何より、博士課程の指導教員Arthur Fineに、デコヒーレンスの分野を研究する自由と指導の両方を与えてくださったことを感謝したい。彼は、当時二人とも初めて接するものだったこの研究分野を、彼と私がよりよく理解できるよう、デコヒーレンスについての「個人的なノート」を書くよう勧めてくれた。そのノートはデコヒーレンスについての総説論文\cite{ref001}に発展し、それが本書の動機になった。この関連で、多くの有益な議論をしてくださり、本書の構想をSpringerの編集者Angela Laheeに伝えてくださったH. Dieter Zehには、深い恩義がある。Angelaはその後、本書の制作のあらゆる側面に、辛抱強く、励ましに満ち、有益な支援を寄せてくださった。

本書の一部を執筆したUniversity of Queenslandでの歓待について、Michael NielsenとGerard Milburnに感謝する。また、第8.4節についてのコメントを寄せてくださったStephen Adler、第3章について意見をくださったErich Joos、量子電気機械系を紹介してくださったGerard Milburn、原稿への多くの貴重なコメントと刺激的な議論をしてくださったWojciech Zurekにも、感謝したい。しかし何にもまして、本書を執筆する長い過程を通じて、あらゆる面で私を励まし、辛抱強く支えてくれた妻Kariに感謝したい。

\begin{flushleft}Melbourne, Australia\hfill Maximilian Schlosshauer\\2007年6月\end{flushleft}
\clearpage
